\chapter{Selected Software and Commonality Analysis}
\label{ch:selected-software-commonality}

This chapter characterizes the final 28 packages (Section~\ref{sec:final-software-set}) for RQ1 and gives the commonality context used later when interpreting the RQ2 measurements. The commonality analysis follows the \citet{Weiss1998} procedure described in Section~\ref{sec:commonality-analysis-procedure}; capabilities are reported only where they could be verified against the project materials.

\section{Overview of Selected Software}
\label{sec:overview-selected-software}

The selected packages come from a mix of development settings---academic research groups, open-source foundations and companies, corporate research labs, and smaller communities or individual maintainers---detailed under V9 (Section~\ref{sec:variabilities}).

The packages span a period of active development from 2001 (ODE's initial release) to the present, with most showing recent repository activity as of the May 2026 measurement date. Licensing, language, and platform details are reported with the measurement data in Section~\ref{sec:repository-project-metadata} (Table~\ref{tab:summary-metadata}).

\section{Functional Software Roles}
\label{sec:robotics-software-categories}

The selected packages can be grouped into three broad functional roles by their primary function in a robotics workflow. These roles are not strict boundaries---some packages span several (see Section~\ref{sec:boundary-cases})---but they structure the commonality analysis.

\subsection{Robotics Simulators}
\label{subsec:robotics-simulators}

Robotics simulators provide virtual environments for representing robots, worlds, sensors, and actuators, letting users test algorithms, validate designs, and collect synthetic data before deploying to hardware. They typically offer graphical rendering, physics simulation, sensor models, and programmatic interfaces for controlling simulated robots. The selected packages primarily classified as simulators are:

\begin{itemize}
    \item \textbf{Gazebo} (including the Ignition/Gazebo Sim lineage) is a long-established open-source 3D robotics simulator developed by Open Robotics. It provides physics simulation through multiple backends, rendering, sensor models, and ROS integration.
    \item \textbf{Webots} is a multi-platform robot simulator maintained by Cyberbotics. It supports modelling, programming, and simulating robots, vehicles, and mechanical systems, with a large library of pre-built robot models.
    \item \textbf{CoppeliaSim} (formerly V-REP) is a versatile robot simulation framework developed by Coppelia Robotics. It supports rapid algorithm development, factory automation simulation, prototyping, verification, and digital twin applications.
    \item \textbf{AirSim} is an open-source simulator for autonomous vehicles (drones and ground vehicles) built on Unreal Engine and developed by Microsoft Research. Its repository-activity classification and officially discontinued status are discussed in Section~\ref{subsec:activity-maintenance-status}.
    \item \textbf{CARLA} is an open-source autonomous driving simulator built on Unreal Engine and developed by the Computer Vision Center in Barcelona and Intel. It provides photorealistic urban environments and a broad sensor suite for autonomous driving research.
    \item \textbf{MORSE} (Modular OpenRobots Simulation Engine) is a generic academic robotic simulator based on Blender and the Bullet physics engine, designed to remain middleware-agnostic across ROS, YARP, and other frameworks.
    \item \textbf{Habitat} is a 3D simulator for embodied AI research developed by Meta AI Research. It supports training and evaluation of embodied AI agents in photorealistic indoor environments.
    \item \textbf{Flightmare} is a flexible modular quadrotor simulator developed by the Robotics and Perception Group at the University of Zurich and ETH Zurich. It separates rendering from the RL environment for flexible configuration between fast and photorealistic modes.
\end{itemize}

Drake, MRPT, and SOFA also provide simulation capabilities as part of broader roles (Section~\ref{sec:boundary-cases}).

\subsection{Physics Engines and Dynamics Libraries}
\label{subsec:physics-engines-dynamics-libraries}

Physics engines and dynamics libraries provide the computational foundation for simulating rigid-body and multibody dynamics, collision detection, and contact resolution. They often serve as backends to higher-level simulators but can also be used directly. The selected packages primarily classified as physics engines or dynamics libraries are:

\begin{itemize}
    \item \textbf{Bullet/PyBullet} is a widely used open-source physics engine for simulation, robotics, and deep reinforcement learning. PyBullet provides Python bindings that have made it popular in the machine learning community.
    \item \textbf{MuJoCo} (Multi-Joint dynamics with Contact) is a general-purpose physics simulator maintained by Google DeepMind. It is designed for fast and accurate simulation of articulated structures and is widely used in robotics and reinforcement learning research.
    \item \textbf{ODE} (Open Dynamics Engine) is one of the earliest open-source libraries for simulating rigid body dynamics. It has been used extensively in robotics and served as a physics backend for Gazebo Classic.
    \item \textbf{PhysX} is NVIDIA's real-time physics engine SDK. The open-source SDK provides capabilities for simulating rigid bodies, particles, vehicles, and deformable bodies, primarily targeted at real-time simulation applications.
    \item \textbf{DART} (Dynamic Animation and Robotics Toolkit) is a C++ physics engine and kinematics/dynamics library developed at Georgia Tech and other institutions. It is designed for stable simulation in manipulation and locomotion research.
    \item \textbf{Simbody} is a multibody dynamics library developed at Stanford University for simulating articulated biomechanical and mechanical systems. It is the physics engine underlying OpenSim (biomedical simulation software).
    \item \textbf{Pinocchio} is a fast and flexible C++ library for rigid-body dynamics computations and their analytical derivatives, developed primarily at Inria and LAAS-CNRS. It implements state-of-the-art algorithms for kinematics, dynamics, and Jacobians.
    \item \textbf{Project Chrono} is an open-source multi-physics simulation library developed at the University of Wisconsin-Madison and the University of Parma. It supports multibody dynamics, finite element analysis, vehicle dynamics, and GPU-accelerated simulation.
    \item \textbf{Newton Dynamics} is a cross-platform physics simulation engine optimized for real-time deterministic rigid-body simulation. It is used in game development and robotics simulation.
    \item \textbf{SOFA} (Simulation Open Framework Architecture) is an open-source real-time simulation framework developed primarily at Inria. It focuses on deformable object simulation, finite element methods, and haptic feedback, with applications in surgical simulation and soft robotics.
\end{itemize}

\subsection{Motion Planning, Middleware, and Manipulation Packages}
\label{subsec:motion-planning-middleware-manipulation}

Motion planning libraries, robotics middleware, and manipulation packages cover parts of the workflow outside full-world simulation: planning algorithms, communication layers, grasp planning, and system integration. They are included because MPS workflows typically combine planning, simulation, and execution rather than using simulators alone. The selected packages primarily classified as planning or middleware packages are:

\begin{itemize}
    \item \textbf{ROS~2} (Robot Operating System 2) is an open-source software development kit and middleware framework for building robot applications. It provides communication infrastructure (based on DDS), hardware abstraction, drivers, libraries, visualization tools, and a large ecosystem of community packages. ROS~2 is the successor to ROS~1 and represents a major redesign with emphasis on real-time control, multi-robot deployment, and production use.
    \item \textbf{MoveIt!} refers to the active MoveIt~2 generation for ROS~2 in this report. It is an open-source robotics manipulation framework that provides motion planning, kinematics, collision checking, trajectory execution, and perception integration for robot arms and manipulators, integrating OMPL and other planners as plugins.
    \item \textbf{OMPL} (Open Motion Planning Library) is an open-source C++ library for sampling-based motion planning algorithms developed at Rice University. It provides implementations of widely used planners such as PRM, RRT, RRT*, and EST, and is the default planning backend for MoveIt!.
    \item \textbf{OpenRAVE} (Open Robotics Automation Virtual Environment) is an open-source robotics planning framework focused on simulation and analysis of kinematic and geometric information for motion planning, particularly for manipulation tasks.
    \item \textbf{Drake} is an open-source C++ and Python toolbox for model-based design and verification of robotics systems, developed by MIT and the Toyota Research Institute (TRI). It covers multibody dynamics, trajectory optimization, model predictive control, perception, and manipulation planning.
    \item \textbf{OROCOS} (Open RObot COntrol Software) is an open-source C++ framework for real-time robot control developed at KU Leuven. Its major components include the Real-Time Toolkit for component-based hard real-time software and the Kinematics and Dynamics Library.
    \item \textbf{MRPT} (Mobile Robot Programming Toolkit) is an open-source cross-platform C++ library for mobile robotics research, providing algorithms for SLAM, computer vision, motion planning, and sensor data processing.
    \item \textbf{YARP} (Yet Another Robot Platform) is open-source middleware for humanoid and general robotics, developed at the Istituto Italiano di Tecnologia. It provides communication infrastructure for distributed robot systems and is the primary middleware for the iCub humanoid robot platform.
    \item \textbf{OpenRTM-\allowbreak aist} implements the OMG Robot Technology Component standard and was developed at AIST in Japan. It provides a component-based framework for creating and connecting robot software modules.
    \item \textbf{GraspIt!} is a simulation environment for robotic grasping research developed at Columbia University. It supports loading robot hand models and 3D object models to plan, evaluate, and visualize grasps.
\end{itemize}

\section{Family and Sub-family Commonalities}
\label{sec:commonalities}

For this analysis, a commonality must hold at the level where it is claimed. A family-wide commonality must hold for every member of the 28-package study set; a feature shared only by one role group is stated as a sub-family commonality; and a feature whose values differ across packages is treated as a variability. This layered structure follows the role grouping in Section~\ref{sec:robotics-software-categories} and avoids presenting a simulator property as if it also characterized middleware or a focused planning library. Figure~\ref{fig:feature-model} summarizes the resulting structure.

\begin{figure}[htbp]
\centering
\begin{tikzpicture}[
  font=\footnotesize,
  root/.style={draw, rounded corners, fill=blue!12, font=\small\bfseries, inner sep=6pt, align=center, text width=9.2cm},
  sub/.style={draw, rounded corners, fill=blue!4, align=left, inner sep=5pt, text width=4.0cm},
  var/.style={draw, rounded corners, fill=orange!7, align=left, inner sep=5pt, text width=12.8cm},
  arr/.style={-{Stealth[length=2mm]}, thick},
]
\node[root] (root) {MPS software family (28 packages)\\[-1pt]
\mdseries C1: MPS-related artifact \quad C2: 3D robot/world state \quad C3: public measurable artifact};
\node[sub, below=11mm of root, xshift=-4.7cm] (s1) {\textbf{S1: Robotics simulators} (8)\\
3D environments and simulation or environment-stepping workflows\\[2pt]
\textit{Associated:} V2, V3, V4, V7};
\node[sub, below=11mm of root] (s2) {\textbf{S2: Physics/dynamics} (10)\\
Physical, kinematic, dynamic, contact, or multiphysics computation\\[2pt]
\textit{Associated:} V3, V4, V5, V6};
\node[sub, below=11mm of root, xshift=4.7cm] (s3) {\textbf{S3: Planning/integration} (10)\\
Planning, manipulation, middleware, and control abstractions\\[2pt]
\textit{Associated:} V1, V3, V4, V7};
\node[var, below=13mm of s2] (vars) {\textbf{Family variability catalogue:}
V1 role; V2 interaction; V3 scope; V4 computation/execution; V5 platform; V6 languages; V7 ecosystem integration; V8 license; V9 governance; V10 repository host.\\[2pt]
V8--V10 cut across all three sub-families; the arrows above identify the role-dependent variability dimensions most directly associated with each sub-family.};
\draw[arr] (root.south) -- (s1.north);
\draw[arr] (root.south) -- (s2.north);
\draw[arr] (root.south) -- (s3.north);
\draw[arr] (s1.south) -- (vars.north west);
\draw[arr] (s2.south) -- (vars.north);
\draw[arr] (s3.south) -- (vars.north east);
\end{tikzpicture}
\caption{Relationship between the MPS family, its three role-based sub-families, and the ten variability dimensions. The sub-family boxes identify their most direct role-dependent variabilities; V8--V10 apply across all three groups.}
\label{fig:feature-model}
\end{figure}

\textbf{C1: MPS-related software artifact.} Every selected package performs at least one task in an MPS workflow: planning, simulation, dynamics, control, manipulation, or integration. This is the top-level commonality. The packages are not interchangeable, but each belongs to the domain this thesis studies, and the sub-families below (with V1--V4) record the different roles they play.

\textbf{C2: Three-dimensional robot or world state.} Every selected package supports three-dimensional robot or environment state. Simulators and physics libraries provide 3D geometry, kinematics, and dynamics directly; planning and manipulation software plans in 3D configuration spaces or workspaces, even when, as in OMPL, the underlying algorithms are dimension-agnostic; middleware carries 3D pose and transform data. How this state is used differs (V3, V4), but support for three dimensions is universal. The one two-dimensional candidate considered, Player/Stage, was not selected (Appendix~\ref{appendix:full-candidate-software-list}).

\textbf{C3: Public measurable technical artifact.} Every package has a public source repository and enough user-facing documentation to be measured. The measurement records confirm a documented API, installation instructions, a getting-started tutorial, and a user manual for all 28 packages. This commonality is partly a product of the selection criteria (Section~\ref{subsec:repository-based-measurability}); it does not mean that all MPS software is equally open or equally documented. License (V8), governance (V9), and repository host (V10) remain variable.

The three sub-families in Table~\ref{tab:subfamily-commonalities} make the role-specific commonalities explicit. These commonalities should not be read as claims about all 28 packages; they hold within the stated sub-family and become variabilities when the whole family is compared.

\begin{table}[H]
\centering
\caption{Sub-family commonalities and associated variabilities.}
\label{tab:subfamily-commonalities}
\singlespacing
\footnotesize
\setlength{\tabcolsep}{2.5pt}
\begin{tabular}{@{}P{2.4cm}P{3.8cm}P{4.3cm}P{2.4cm}@{}}
\toprule
\textbf{Sub-family} & \textbf{Members} & \textbf{Sub-family commonalities} & \textbf{Associated variabilities} \\
\midrule
S1: Robotics simulators & AirSim, CARLA, CoppeliaSim, Flightmare, Gazebo, Habitat, MORSE, Webots & Provide 3D virtual environments for representing robots, worlds, sensors, or agents; support a simulation or environment-stepping workflow used for testing, data generation, or algorithm development. & Interaction mode (V2), application scope (V3), computation/execution model (V4), ecosystem integration (V7) \\
\hline
S2: Physics engines and dynamics libraries & Bullet/PyBullet, DART, MuJoCo, Newton Dynamics, ODE, PhysX, Pinocchio, Project Chrono, Simbody, SOFA & Compute physical, kinematic, or dynamic properties of bodies, articulated systems, contacts, or deformable/multiphysics models; expose solver or library interfaces for use directly or as backends. & Application scope (V3), computation model (V4), platform support (V5), programming languages (V6) \\
\hline
S3: Planning, middleware, and manipulation software & Drake, GraspIt!, MoveIt!, MRPT, OMPL, OpenRAVE, OpenRTM-\allowbreak aist, OROCOS, ROS~2, YARP & Expose reusable application-level robotics abstractions rather than standalone physics backends. Planning and manipulation packages express these abstractions as motions and grasps; middleware and control packages express them as components, devices, controllers, and messages. & Primary role (V1), application scope (V3), computation/execution model (V4), ecosystem integration (V7) \\
\bottomrule
\end{tabular}
\end{table}

\section{Variabilities}
\label{sec:variabilities}

The family-wide and sub-family commonalities above are realized differently across the selected packages. The ten variability dimensions below include role-specific patterns such as physics or kinematics computation, ecosystem integration, and simulation or planning loops. These patterns are common within some sub-families but are variabilities at the level of the full 28-package family.

\textbf{V1: Primary Software Role.} The most fundamental variability is the role a package plays in a robotics workflow, ranging from full simulation environments to physics engines to planning libraries to middleware frameworks (Section~\ref{sec:robotics-software-categories}).

\textbf{V2: Interaction Mode.} Some packages combine a graphical environment with a programmatic API (Gazebo, Webots, and CoppeliaSim), whereas focused libraries are primarily API-driven (OMPL, Pinocchio, and ODE). This variability describes the user-facing interaction mode, not the C3 measurement baseline that all selected packages provide enough technical material to be studied.

\textbf{V3: Application Scope.} General-purpose robotics simulators and engines coexist with packages specialized for autonomous driving, aerial vehicles, embodied AI, medical or soft robotics, grasping, mobile robotics, manipulation, and real-time control (Section~\ref{sec:robotics-software-categories}).

\textbf{V4: Computation or Execution Model.} The software may implement a time-stepped simulation loop, a dynamics or kinematics solver, a planning or optimization pipeline, a middleware/component runtime, or several of these models in one toolbox. This variability captures capabilities that are common within some sub-families, such as physics computation and iterative simulation or planning loops, but not across all 28 packages.

\textbf{V5: Platform Support.} All packages support Linux, and support for other platforms varies. Some packages provide first-class support for Windows and macOS (Webots, CoppeliaSim, Bullet), while others are primarily developed for Linux with limited cross-platform support.

\textbf{V6: Recorded Programming Languages.} C++ and Python are the dominant languages in the measurement records, but the recorded sets differ. Some packages use C or C++ without Python, many pair C++ with Python, and larger frameworks record additional languages such as Java, MATLAB, or IDL. This variability reports the languages in the measurement dataset; it does not infer a public API from repository language counts.

\textbf{V7: Robotics-Ecosystem Integration.} MoveIt! is ROS-native, Gazebo is closely ROS-integrated, and simulators such as Webots and CARLA provide bridges. MORSE supports several middleware systems. YARP, OROCOS, and OpenRTM-\allowbreak aist provide alternative middleware or component models, while other libraries are usable independently. Ecosystem integration is therefore a variability with native, bridge, multi-middleware, alternative-middleware, and standalone values.

\textbf{V8: License Model.} The measured source artifacts use permissive, reciprocal, dual, and mixed distribution models. CoppeliaSim is the principal mixed-model case because the measured GPL source coexists with commercial and educational distribution licences. License choice affects how software can be redistributed and integrated.

\textbf{V9: Governance Model.} Governance ranges across companies, open-source foundations, academic or public research institutions, mixed consortia, and independent communities or maintainers. The categories describe the primary stewardship visible in the measurement snapshot rather than legal ownership of every component.

\textbf{V10: Repository Host.} GitHub hosts the measured source for 27 packages. ODE is the exception: its primary repository is on Bitbucket, with a GitHub mirror also recorded. The commonality is public Git-based measurability (C3), while the hosting service is variable.

\section{Parameters of Variation}
\label{sec:parameters-of-variation}

For each variability identified in Section~\ref{sec:variabilities}, Table~\ref{tab:parameters-of-variation} summarizes the parameters of variation and their typical binding times (defined in Section~\ref{sec:commonality-analysis-background}). Binding time records when a value is fixed for a package or selected by an integrator; it is not the date on which the measurement was collected.

\begin{table}[H]
\centering
\caption{Parameters of variation and binding times.}
\label{tab:parameters-of-variation}
\singlespacing
\small
\begin{tabular}{@{}P{3.2cm}P{7.1cm}P{3.0cm}@{}}
\toprule
\textbf{Variability} & \textbf{Parameters of Variation} & \textbf{Binding Time} \\
\midrule
V1: Primary role & Simulator; physics/dynamics engine, library, or framework; planning or manipulation library/framework; middleware or control framework; multi-purpose or multiphysics package & Specification time \\
\hline
V2: Interaction mode & Graphical plus API, primarily API & Specification / run time \\
\hline
V3: Application scope & General robotics/physics; autonomous vehicles; aerial vehicles; embodied AI; medical/soft robotics; grasping; mobile robotics; motion/manipulation planning; real-time control; rigid-body, multibody, or multiphysics applications & Specification time \\
\hline
V4: Computation/execution & Simulation loop; dynamics or kinematics solver; planning, search, or optimization pipeline; middleware/component runtime; combined planning/simulation, dynamics/simulation, component/kinematics, or multi-role model & Specification / configuration time \\
\hline
V5: Platform support & Linux, with the measured combinations of Windows, macOS, Android, or BSD support & Specification / build time \\
\hline
V6: Recorded languages & C++; C/C++; C/Python; C++/Python; multi-language; unclear in the measurement record & Build / configuration time \\
\hline
V7: Ecosystem integration & ROS-native/integrated, ROS bridge, multi-middleware, alternative middleware/component model, standalone & Build / configuration / run time \\
\hline
V8: License & Specific permissive, reciprocal, dual, or mixed open-source/commercial licence recorded in the sub-family package maps below & Specification time \\
\hline
V9: Governance & Company, foundation, academic/public institution, mixed consortium, independent/community & Specification time \\
\hline
V10: Repository host & GitHub, Bitbucket with GitHub mirror & Specification time \\
\bottomrule
\end{tabular}
\end{table}

The binding times reflect when a user or integrator can still make a choice. Roles, governance, and repository hosting are fixed at specification or project-organization time, whereas ecosystem integration may also be selected when a package is built, configured, or connected at run time.

Tables~\ref{tab:simulator-variability-map}--\ref{tab:planning-variability-map} map all ten variabilities at package level, one table per sub-family. This arrangement keeps a package's functional, delivery, and governance values in one row. The values synthesize the final-set roles in Table~\ref{tab:final-software-set}, the metadata in Table~\ref{tab:summary-metadata}, and the package purposes in the measurement dataset; they are descriptive, not quality scores.

For compactness, L, W, M, A, and B denote Linux, Windows, macOS, Android, and BSD; GH and BB denote GitHub and Bitbucket; and MW denotes middleware. V7 distinguishes ROS-native or closely integrated software, optional ROS bridges, multi-middleware support, alternative middleware/component models, and standalone use.

\newcommand{\variationmapheader}{%
\toprule
\textbf{Software} & \textbf{V1 Role} & \textbf{V2 Interact.} & \textbf{V3 Scope} & \textbf{V4 Execution} & \textbf{V5 Platform} & \textbf{V6 Languages} & \textbf{V7 Integration} & \textbf{V8 License} & \textbf{V9 Governance} & \textbf{V10 Host} \\
\midrule}

\begin{landscape}
\begin{table}[H]
\centering
\caption{Complete V1--V10 package map for S1: robotics simulators.}
\label{tab:simulator-variability-map}
\begingroup
\singlespacing
\fontsize{7pt}{8.2pt}\selectfont
\setlength{\tabcolsep}{1.2pt}
\begin{tabular}{@{}P{2.0cm}P{2.2cm}P{1.25cm}P{2.15cm}P{2.35cm}P{1.35cm}P{1.65cm}P{1.65cm}P{1.4cm}P{1.65cm}P{1.05cm}@{}}
\variationmapheader
AirSim & Simulator & GUI + API & Autonomous vehicles & Simulation loop & L/W & C++/\allowbreak{}Python & ROS bridge & MIT & Company & GH \\
CARLA & Simulator & GUI + API & Autonomous driving & Simulation loop & L/W & C++/\allowbreak{}Python & ROS bridge & MIT & Mixed consortium & GH \\
CoppeliaSim & Simulator & GUI + API & General robotics & Simulation loop & L/W/M & C++/\allowbreak{}Python & ROS bridge & GPL + commercial & Company & GH \\
Flightmare & Simulator & GUI + API & Aerial vehicles & Simulation loop & L & C++/\allowbreak{}Python & Standalone & MIT & Academic & GH \\
Gazebo & Simulator & GUI + API & General robotics & Simulation loop & L/W/M & C++/\allowbreak{}Python & ROS integrated & Apache-2.0 & Foundation & GH \\
Habitat & Simulator & API & Embodied AI & Simulation loop & L/M & C++/\allowbreak{}Python & Standalone & MIT & Company & GH \\
MORSE & Simulator & GUI + API & General robotics & Simulation loop & L/W/M & C/\allowbreak{}Python & Multi-MW & BSD & Academic & GH \\
Webots & Simulator & GUI + API & General robotics & Simulation loop & L/W/M & Multi-language & ROS bridge & Apache-2.0 & Company & GH \\
\bottomrule
\end{tabular}
\endgroup
\end{table}
\end{landscape}

\begin{landscape}
\begin{table}[H]
\centering
\caption{Complete V1--V10 package map for S2: physics engines and dynamics libraries.}
\label{tab:physics-variability-map}
\begingroup
\singlespacing
\fontsize{7pt}{8.2pt}\selectfont
\setlength{\tabcolsep}{1.2pt}
\begin{tabular}{@{}P{2.0cm}P{2.2cm}P{1.25cm}P{2.15cm}P{2.35cm}P{1.35cm}P{1.65cm}P{1.65cm}P{1.4cm}P{1.65cm}P{1.05cm}@{}}
\variationmapheader
Bullet/PyBullet & Physics engine & API & General physics/robotics & Dynamics solver & L/W/M & C++/\allowbreak{}Python & Standalone & zlib & Community & GH \\
DART & Dynamics library & API & General robotics/\allowbreak{}dynamics & Dynamics/kinematics solver & L/W/M/\allowbreak{}B & C++/\allowbreak{}Python & Standalone & BSD & Academic & GH \\
MuJoCo & Physics engine & GUI + API & General physics/robotics & Dynamics solver & L/W/M & C++/\allowbreak{}Python & Standalone & Apache-2.0 & Company & GH \\
Newton Dynamics & Physics engine & API & General physics & Dynamics solver & L/W/M/\allowbreak{}A & C/\allowbreak{}C++ & Standalone & zlib & Independent & GH \\
ODE & Physics engine & API & General physics & Dynamics solver & L/W/M & Unclear & Standalone & BSD/\allowbreak{}LGPL & Independent & BB + GH \\
PhysX & Physics engine & API & General physics & Dynamics solver & L/W & C++ & Standalone & BSD & Company & GH \\
Pinocchio & Dynamics library & API & Rigid-body robotics & Dynamics/kinematics solver & L/W/M & C++/\allowbreak{}Python & Standalone & BSD & Academic & GH \\
Project Chrono & Multi-physics simulation & API & Multiphysics/vehicles & Dynamics/simulation & L/W/M & C++/\allowbreak{}Python & Standalone & BSD & Academic & GH \\
Simbody & Dynamics library & API & Multibody/\allowbreak{}biomechanics & Dynamics solver & L/W/M & C/\allowbreak{}C++ & Standalone & Apache-2.0 & Academic & GH \\
SOFA & Simulation framework & GUI + API & Medical/soft robotics & Simulation loop & L/W/M & C++/\allowbreak{}Python & Standalone & LGPL-2.1 & Academic & GH \\
\bottomrule
\end{tabular}
\endgroup
\end{table}
\end{landscape}

\begin{landscape}
\begin{table}[H]
\centering
\caption{Complete V1--V10 package map for S3: planning, middleware, and manipulation software.}
\label{tab:planning-variability-map}
\begingroup
\singlespacing
\fontsize{7pt}{8.2pt}\selectfont
\setlength{\tabcolsep}{1.2pt}
\begin{tabular}{@{}P{2.0cm}P{2.2cm}P{1.25cm}P{2.15cm}P{2.35cm}P{1.35cm}P{1.65cm}P{1.65cm}P{1.4cm}P{1.65cm}P{1.05cm}@{}}
\variationmapheader
Drake & Multi-purpose toolbox & API & General robotics & Multi-role & L/M & C++/\allowbreak{}Python & Standalone & BSD & Mixed consortium & GH \\
GraspIt! & Manipulation framework & GUI + API & Robotic grasping & Planning/simulation & L/W/M & C/\allowbreak{}C++/\allowbreak{}IDL & Standalone & GPL v3+ & Academic & GH \\
MoveIt! & Manipulation framework & GUI + API & Manipulation & Planning pipeline & L & C++/\allowbreak{}Python & ROS native & BSD & Mixed consortium & GH \\
MRPT & Multi-purpose library & GUI + API & Mobile robotics & Multi-role & L/W/M & C++/\allowbreak{}Python & ROS bridge & BSD & Academic & GH \\
OMPL & Motion planning library & API & Motion planning & Planning/search & L/W/M & C++/\allowbreak{}Python & Standalone & BSD & Academic & GH \\
OpenRAVE & Planning framework & GUI + API & Manipulation/planning & Planning pipeline & L/W/M & C++/\allowbreak{}Python & ROS bridge & LGPL/\allowbreak{}A2.0 & Mixed & GH \\
OpenRTM-\allowbreak aist & Middleware & GUI + API & General robotics & Component runtime & L/W/M & C++/\allowbreak{}Python & Alternative MW & LGPL-2.1 & Public institute & GH \\
OROCOS & Control framework & API & Real-time control & Component/kinematics & L/W/M & C++/\allowbreak{}Python & Alternative MW & GPL/\allowbreak{}LGPL & Academic/\allowbreak{}community & GH \\
ROS~2 & Middleware and SDK & GUI + API & General robotics & Middleware runtime & L/W/M & Multi-language & ROS native & Apache-2.0 & Foundation & GH \\
YARP & Middleware & API & Humanoid/general robotics & Middleware runtime & L/W/M & Multi-language & Alternative MW & BSD & Public institute & GH \\
\bottomrule
\end{tabular}
\endgroup
\end{table}
\end{landscape}

\let\variationmapheader\relax

\section{Boundary Cases}
\label{sec:boundary-cases}

Several of the selected packages occupy intermediate positions in the classification above; the boundaries between roles are not always sharp.

\textbf{ROS~2} is classified as middleware in this analysis, but its scope extends well beyond communication infrastructure. It is the de facto integration layer for robotics software components, and many other packages in the selected set assume it as their runtime environment.

\textbf{Drake} spans multiple roles. It provides simulation capabilities without being a full simulator in the sense of Gazebo or Webots. It includes sophisticated multibody dynamics without being solely a physics engine. It offers trajectory optimization and model predictive control without being only a planning library.

The boundary between \textbf{MoveIt!} and \textbf{OMPL} reflects architectural layering rather than functional separation: OMPL is a general-purpose planning library usable independently, while MoveIt! builds kinematics, collision checking, and trajectory execution around it (Section~\ref{subsec:motion-planning-middleware-manipulation}).

\textbf{SOFA} is grouped with the physics engines because the measured materials emphasize deformable-body physics rather than whole-robot simulation (Section~\ref{subsec:physics-engines-dynamics-libraries}); it is a boundary case rather than a direct peer of rigid-body engines such as ODE or Bullet.

\textbf{GraspIt!} sits between the manipulation and simulator roles. It is grouped with the manipulation packages because grasp planning is its purpose, but in practice it runs as a standalone grasp simulator: it loads hand and object models, simulates contacts, and scores candidate grasps through its own interface. Its robot-level grasp abstractions satisfy the sub-family commonality in Table~\ref{tab:subfamily-commonalities}, although users generally run GraspIt! as an application rather than embed it as a component.

In the chapters that follow, measurements are reported for each package individually and interpreted in the context of its role, so these boundary cases are not treated as measurement anomalies.
